\section{Related Work}
\label{sec:related}

We survey five bodies of work that Omni-Shield builds on
or directly compares against: PII detection systems,
vision-language models for document understanding,
knowledge distillation for specialised domains,
cryptographic audit mechanisms, and multi-agent
orchestration frameworks.

% ─────────────────────────────────────────────────────────
\subsection{PII Detection and Document Redaction}
% ─────────────────────────────────────────────────────────

The dominant paradigm for automated PII detection is
named entity recognition (NER) applied to extracted text.
Microsoft Presidio~\cite{presidio2020} combines spaCy
NER with regex pattern matching to identify entities
such as names, phone numbers, and national identifiers.
Presidio is open-source, language-extensible, and widely
deployed, but operates exclusively on text: it cannot
process image, audio, or video inputs, and has no
mechanism for detecting visual PII such as signatures
or face photographs.
Amazon Comprehend~\cite{aws_comprehend} and Google Cloud
Data Loss Prevention~\cite{google_dlp} offer cloud-based
PII detection with broader entity coverage and higher
throughput, but transmit documents to third-party
infrastructure --- a fundamental incompatibility with
privacy-sensitive workflows where documents must not
leave an organisation's control.

Specialised document redaction for identity documents
has received comparatively little academic attention.
Prior work focuses on redaction of court
documents~\cite{berg2020}, medical records~\cite{stubbs2015},
and legal contracts~\cite{manor2019}, where text
extraction is straightforward and PII is confined to
running prose.
Identity documents present a structurally different
problem: PII is densely packed, visually structured,
and partially non-textual (faces, signatures, barcodes).
The \textit{DocRedact} system~\cite{ferrando2022} applies
LayoutLM embeddings to redact form-structured documents,
but requires pre-labelled field schemas and does not
generalise to unseen card designs.

In contrast to all prior work, Omni-Shield processes
documents at the pixel level using a VLM that
understands layout, operates without pre-defined schemas,
handles six document modalities, and runs entirely
on-device.

% ─────────────────────────────────────────────────────────
\subsection{Vision-Language Models for Document
            Understanding}
% ─────────────────────────────────────────────────────────

The application of vision-language models to document
understanding has progressed rapidly.
LayoutLM~\cite{xu2020layoutlm} and its successors
LayoutLMv2 and LayoutLMv3~\cite{xu2021layoutlmv2,
huang2022layoutlmv3} demonstrate that jointly encoding
text tokens with their spatial positions on a page
substantially improves document classification, key
information extraction, and form understanding.
Donut~\cite{kim2022donut} removes the explicit OCR
dependency by training an end-to-end encoder-decoder
to read documents directly from pixels.
DocVQA~\cite{mathew2021docvqa} establishes visual
question answering over scanned documents as a
benchmark task.

General-purpose multimodal models such as
GPT-4V~\cite{openai2023gpt4v},
Gemini~\cite{team2023gemini}, and the Qwen-VL
family~\cite{bai2023qwenvl} demonstrate strong
zero-shot document understanding,
but at parameter counts (7B--175B) that preclude
deployment on consumer-grade hardware.
Qwen2-VL-2B-Instruct~\cite{wang2024qwen2vl}, the base
model we fine-tune, is specifically designed for
efficient edge deployment and achieves competitive
document understanding performance at 2 billion
parameters.

Our work differs from prior document VLM research in
two respects.
First, we fine-tune the VLM specifically for PII
\emph{localisation} --- producing bounding box coordinates
alongside entity labels --- rather than for general
question answering.
Second, we embed the VLM as one component of a
multi-agent pipeline rather than using it as a
standalone system, allowing other agents to correct,
filter, and refine its outputs.

% ─────────────────────────────────────────────────────────
\subsection{Knowledge Distillation for Specialised Domains}
% ─────────────────────────────────────────────────────────

Knowledge distillation~\cite{hinton2015distilling}
transfers the generalisation capability of a large
teacher model into a smaller student model by training
the student on teacher-generated soft labels rather
than hard ground truth.
The approach has been applied to compression of BERT
models~\cite{sanh2019distilbert}, structured prediction
tasks~\cite{li2022explaining}, and most recently to
instruct-following for instruction-tuned LLMs using
synthetic data generated by a more capable
model~\cite{taori2023alpaca, wang2023selfinstruct}.

Parameter-efficient fine-tuning methods --- in
particular LoRA~\cite{hu2022lora} and its
4-bit quantised variant QLoRA~\cite{dettmers2023qlora}
--- have made it practical to adapt large models on
consumer hardware without full parameter updates.
QLoRA reduces peak GPU memory during fine-tuning by
introducing double quantisation and paged optimisers,
enabling fine-tuning of 7B+ models on a single 24\,GB
GPU and, in our case, a 2B VLM on an 8\,GB GPU.

Our knowledge distillation methodology follows the
teacher-student paradigm with synthetic data:
we generate 2,000 synthetic identity documents and
label them using the Claude API as teacher~\cite{anthropic2024},
then train the student VLM on these structured
annotations.
A key finding is that the primary contribution of
fine-tuning is \emph{output format adherence} rather
than detection accuracy --- the adapter trains the
model to emit consistently parseable
\texttt{LABEL\,|\,VALUE\,|\,BBOX} triples, which the
downstream pipeline requires for reliable extraction.
This is consistent with observations in instruction
fine-tuning literature that format alignment is a
distinct and learnable skill separate from factual
knowledge~\cite{ouyang2022instructgpt}.

% ─────────────────────────────────────────────────────────
\subsection{Blockchain Audit Trails for Data Processing}
% ─────────────────────────────────────────────────────────

Blockchain technology has been applied to data
provenance and audit in healthcare~\cite{xia2017medshare},
supply chain~\cite{Casino2019}, and
legal evidence preservation~\cite{zhu2019blockchain}.
The core property that makes blockchain suitable for
audit is tamper-evidence: once a hash is anchored
on-chain, it cannot be altered retroactively without
invalidating the entire subsequent chain.

Smart contract platforms such as
Ethereum~\cite{buterin2014ethereum} allow the audit
logic itself (who may record, what is recorded, how
records are retrieved) to be encoded in verifiable
on-chain code.
Gas cost is a practical constraint: complex audit
contracts on Ethereum mainnet can be prohibitively
expensive, motivating the use of Layer~2 solutions such
as Polygon~\cite{polygon2021} for production deployments.

Prior work on blockchain-based document audit focuses
primarily on \emph{integrity} --- proving that a
document has not been altered since it was recorded.
Omni-Shield extends this to \emph{redaction
provenance}: proving not just that the document exists
but that a specific redaction operation was performed
on it, by a specific owner, at a specific time.
This distinction is important for GDPR Article~5(2)
accountability obligations, which require demonstrating
that personal data was processed in compliance with
the regulation --- a much stronger claim than mere
integrity.

% ─────────────────────────────────────────────────────────
\subsection{Zero-Knowledge Proofs for Privacy}
% ─────────────────────────────────────────────────────────

Zero-knowledge proofs (ZKPs) allow a prover to
convince a verifier that a statement is true without
revealing any information beyond the truth of the
statement~\cite{goldwasser1989knowledge}.
Succinct non-interactive arguments of knowledge
(SNARKs)~\cite{groth2016size} make this practical for
real deployments: the Groth16 proving system produces
constant-size proofs (128 bytes) verifiable in
milliseconds regardless of the complexity of the
proved statement.

ZoKrates~\cite{eberhardt2018zokrates} provides a
high-level language and toolchain for compiling
programs into Groth16/BN128 circuits, lowering the
barrier to SNARK implementation significantly.
ZKPs have been applied to private credential
presentation~\cite{camenisch1997efficient},
anonymous voting~\cite{cortier2019belenios},
and confidential blockchain
transactions~\cite{sasson2014zerocash}.

To our knowledge, Omni-Shield is the first system to
apply ZK-SNARKs specifically to document redaction
provenance: the proof attests that a document with
a known hash was redacted, and that the redaction
committed to a specific (undisclosed) set of PII
regions, without revealing which regions were redacted
or what values they contained.
This provides a \emph{selective disclosure} property
particularly relevant for legal and compliance contexts:
an auditor can verify that redaction occurred without
gaining access to the redacted content.

% ─────────────────────────────────────────────────────────
\subsection{Multi-Agent Systems for Document Processing}
% ─────────────────────────────────────────────────────────

Multi-agent LLM systems decompose complex tasks into
specialised sub-tasks handled by coordinated agents.
AutoGen~\cite{wu2023autogen} and
MetaGPT~\cite{hong2023metagpt} demonstrate that
agent pipelines can outperform monolithic LLMs on
tasks requiring planning, tool use, and iterative
refinement.
In document processing, agent decomposition allows
different agents to specialise: one for layout
understanding, one for PII classification, one for
bounding box refinement, one for false-positive
filtering.

The CriticAgent pattern --- a dedicated agent that
reviews and challenges the output of prior agents ---
mirrors the adversarial review step in self-refine
frameworks~\cite{madaan2023selfrefine} and reduces
false-positive rates in our pipeline.
Our ablation study quantifies the contribution of
each agent class, finding that visual detection
agents (LayoutAgent via VLM, face detection via
YOLOv8~\cite{redmon2016yolo}) account for 54\% of
total F1 --- substantially more than text-based
agents, validating the multi-modal decomposition.

\vspace{6pt}
\noindent\textbf{Positioning.}
Omni-Shield occupies a distinct position in this
landscape: it is the only system that simultaneously
achieves local-first deployment, multi-modal input
(including visual PII), and cryptographically
verifiable redaction audit on consumer hardware.
Table~\ref{tab:comparison} in Section~\ref{sec:evaluation}
quantifies these differences against the closest
alternatives.
